%% ============================================================
%% Policy Gradient Methods for Continuous Decision Processes
%% Self Project Report — July 2025
%% ============================================================

\documentclass[11pt, a4paper]{article}

% ── Packages ──────────────────────────────────────────────────
\usepackage[margin=2.5cm]{geometry}
\usepackage{amsmath, amssymb, amsthm, mathtools}
\usepackage{bm}           % Bold math
\usepackage{bbm}          % Indicator 1
\usepackage[T1]{fontenc}
\usepackage[utf8]{inputenc}
\usepackage{lmodern}
\usepackage{microtype}
\usepackage{hyperref}
\usepackage{xcolor}
\usepackage{graphicx}
\usepackage{subcaption}
\usepackage{booktabs}
\usepackage{array}
\usepackage{multirow}
\usepackage{algorithm}
\usepackage{algpseudocode}
\usepackage{listings}
\usepackage{tikz}
\usepackage{pgfplots}
\pgfplotsset{compat=1.18}
\usepackage{enumitem}
\usepackage{tcolorbox}
\usepackage[backend=bibtex, style=authoryear, natbib=true]{biblatex}
\usepackage{cleveref}

% ── Colours ───────────────────────────────────────────────────
\definecolor{deepblue}{RGB}{30,100,170}
\definecolor{codegreen}{RGB}{40,140,60}
\definecolor{codegray}{RGB}{128,128,128}
\definecolor{codebg}{RGB}{248,248,248}
\definecolor{theorembg}{RGB}{240,248,255}
\definecolor{defbg}{RGB}{255,250,240}
\definecolor{algbg}{RGB}{248,255,248}

% ── Hyperref ──────────────────────────────────────────────────
\hypersetup{
    colorlinks=true,
    linkcolor=deepblue,
    citecolor=codegreen,
    urlcolor=deepblue,
}

% ── Theorem environments ──────────────────────────────────────
\tcbuselibrary{theorems, skins}
\newtcbtheorem[number within=section]{theorem}{Theorem}{
    colback=theorembg, colframe=deepblue!70,
    fonttitle=\bfseries, separator sign none,
    description delimiters parenthesis,
}{thm}
\newtcbtheorem[number within=section]{definition}{Definition}{
    colback=defbg, colframe=orange!60,
    fonttitle=\bfseries, separator sign none,
}{def}
\newtcbtheorem[number within=section]{proposition}{Proposition}{
    colback=algbg, colframe=codegreen!70,
    fonttitle=\bfseries, separator sign none,
}{prop}
\newtcbtheorem[number within=section]{remark}{Remark}{
    colback=gray!8, colframe=gray!50,
    fonttitle=\bfseries, separator sign none,
}{rem}

% ── Code listings ─────────────────────────────────────────────
\lstset{
    backgroundcolor=\color{codebg},
    commentstyle=\color{codegray}\itshape,
    keywordstyle=\color{deepblue}\bfseries,
    stringstyle=\color{codegreen},
    basicstyle=\ttfamily\small,
    breakatwhitespace=false, breaklines=true,
    captionpos=b, keepspaces=true,
    numbers=left, numbersep=5pt,
    numberstyle=\tiny\color{codegray},
    showspaces=false, showstringspaces=false,
    showtabs=false, tabsize=2,
    frame=single, rulecolor=\color{gray!40},
}
\lstdefinestyle{Python}{language=Python}

% ── Maths macros ──────────────────────────────────────────────
\newcommand{\R}{\mathbb{R}}
\newcommand{\E}{\mathbb{E}}
\newcommand{\N}{\mathcal{N}}
\newcommand{\Prob}{\mathbb{P}}
\newcommand{\st}{\mathcal{S}}
\newcommand{\ac}{\mathcal{A}}
\newcommand{\tr}{\mathrm{tr}}
\newcommand{\diag}{\mathrm{diag}}
\newcommand{\T}{\top}
\newcommand{\grad}{\nabla}
\newcommand{\hgrad}{\hat{\nabla}}
\newcommand{\policy}{\pi_{\theta}}
\newcommand{\traj}{\tau}
\newcommand{\KL}{\mathrm{KL}}
\newcommand{\FIM}{\mathbf{F}}
\newcommand{\norm}[1]{\left\lVert #1 \right\rVert}

% ── Bibliography ──────────────────────────────────────────────
\addbibresource{references.bib}

% ── Title ─────────────────────────────────────────────────────
\title{
    \vspace{-1cm}
    {\Large\bfseries Policy Gradient Methods for Continuous Decision Processes}\\[0.5em]
    {\large Self Project Report}
}
\author{
    Self Project\\
    \small July 2025
}
\date{}

% ==============================================================
\begin{document}
% ==============================================================

\maketitle
\vspace{-0.5cm}
\hrule
\vspace{0.3cm}

\begin{abstract}
We present a comprehensive study of policy gradient reinforcement learning
for continuous control problems formulated as Markov Decision Processes (MDPs).
Starting from first principles, we derive the policy gradient theorem, implement
stochastic Gaussian policies parameterised by deep neural networks, and develop
several gradient estimators including REINFORCE \citep{williams1992}, 
Advantage Actor-Critic (A2C), and the Natural Policy Gradient \citep{kakade2001}.
A central focus is variance reduction: we empirically study the effect of
baseline subtraction on gradient estimator signal-to-noise ratio and show that
neural network value baselines reduce variance by up to 60\% compared to the
no-baseline estimator.
We extend the framework to continuous action spaces via Gaussian policies and
apply it to linear-quadratic regulator (LQR) control and nonlinear pendulum
swing-up tasks.
Finally, we establish and experimentally verify the connection between
policy gradient RL and stochastic optimal control by comparing learned value
functions against the analytic Hamilton--Jacobi--Bellman (HJB) solution obtained
through the Discrete Algebraic Riccati Equation.
Our implementation achieves less than 5\% regret relative to the analytic
optimal LQR controller on a 4-dimensional system, and successfully learns
to swing up and balance a continuous pendulum without access to system dynamics.
\end{abstract}

\tableofcontents
\newpage

% ==============================================================
\section{Introduction}
% ==============================================================

Sequential decision-making under uncertainty is a fundamental challenge
in engineering, robotics, and operations research.
The reinforcement learning (RL) paradigm formalises this as an agent interacting
with an environment modelled as a Markov Decision Process (MDP), seeking to
maximise expected cumulative reward without direct knowledge of system dynamics
or the reward structure's optimum.

\textbf{Motivation.}
Classical optimal control methods --- such as Linear Quadratic Regulation (LQR)
and model predictive control --- require explicit knowledge of system dynamics
$f(x, u)$ and can be computationally intractable for nonlinear or high-dimensional
systems.
Policy gradient RL sidesteps these requirements by directly optimising a
parameterised policy through gradient ascent on expected return, making it
applicable to problems where dynamics are unknown, partially observed, or
too complex for analytic treatment.

\textbf{Continuous action spaces.}
Many physical control problems --- robotics, energy systems, process control ---
have continuous action spaces.
This necessitates policies that output probability distributions over continuous
actions, for which Gaussian policies provide a natural, differentiable, and
analytically tractable choice.

\textbf{Goals of this project.}
This project has four interconnected goals:
\begin{enumerate}[leftmargin=1.5em]
    \item Formulate control problems as MDPs and derive the policy gradient
    theorem from first principles.
    \item Implement and compare policy gradient algorithms (REINFORCE, A2C, NPG)
    with rigorous variance analysis.
    \item Extend the framework to continuous action spaces using Gaussian policies
    and stochastic optimisation.
    \item Establish the connection between learned policies and the analytic
    HJB/DARE optimum for LQR systems.
\end{enumerate}

\textbf{Organisation.}
\Cref{sec:mdp} introduces MDPs and the policy optimisation objective.
\Cref{sec:pg} derives the policy gradient theorem and its practical estimators.
\Cref{sec:variance} analyses variance and baseline subtraction.
\Cref{sec:algorithms} presents the algorithms implemented.
\Cref{sec:gaussian} develops Gaussian policies for continuous action spaces.
\Cref{sec:optimal_control} connects policy gradient RL with stochastic optimal
control and the HJB equation.
\Cref{sec:experiments} presents our experimental results.
\Cref{sec:conclusion} concludes and identifies directions for future work.

% ==============================================================
\section{Markov Decision Processes}
\label{sec:mdp}
% ==============================================================

\begin{definition}{Markov Decision Process}{}
A \emph{Markov Decision Process} (MDP) is a tuple
$\mathcal{M} = (\mathcal{S}, \mathcal{A}, P, R, \gamma)$ where:
\begin{itemize}
    \item $\mathcal{S} \subseteq \R^n$ is the state space;
    \item $\mathcal{A} \subseteq \R^m$ is the (continuous) action space;
    \item $P : \mathcal{S} \times \mathcal{A} \times \mathcal{S} \to \R_{\geq 0}$
    is the transition kernel, with $P(s' \mid s, a)$ the density of the next state;
    \item $R : \mathcal{S} \times \mathcal{A} \to \R$ is the expected reward function;
    \item $\gamma \in [0, 1)$ is the discount factor.
\end{itemize}
\end{definition}

The Markov property requires that the next state $s_{t+1}$ is conditionally
independent of the history $(s_0, a_0, \ldots, s_{t-1}, a_{t-1})$ given $(s_t, a_t)$:
\begin{equation}
    P(s_{t+1} \mid s_0, a_0, \ldots, s_t, a_t) = P(s_{t+1} \mid s_t, a_t).
\end{equation}

\subsection{Policies and Trajectories}

A \emph{stochastic policy} $\pi : \mathcal{S} \times \mathcal{A} \to \R_{\geq 0}$
is a conditional density over actions given states.
A \emph{trajectory} $\tau = (s_0, a_0, s_1, a_1, \ldots, s_T, a_T)$ is generated
by initialising $s_0 \sim \mu_0$, then iterating:
\begin{equation}
    a_t \sim \pi(\cdot \mid s_t), \qquad s_{t+1} \sim P(\cdot \mid s_t, a_t).
\end{equation}

The probability density of a trajectory under $\pi$ factorises as:
\begin{equation}
    p_\pi(\tau) = \mu_0(s_0) \prod_{t=0}^{T} \pi(a_t \mid s_t) \cdot P(s_{t+1} \mid s_t, a_t).
\end{equation}

\subsection{The Policy Optimisation Objective}

We seek a policy $\pi_\theta$ parameterised by $\theta \in \R^d$ that maximises
the \emph{expected discounted return}:
\begin{equation}
    J(\theta) = \E_{\tau \sim p_{\pi_\theta}} \left[ G(\tau) \right], \qquad
    G(\tau) = \sum_{t=0}^{T} \gamma^t \, r(s_t, a_t).
    \label{eq:objective}
\end{equation}

The \emph{state value function} and \emph{action-value function} are:
\begin{align}
    V^{\pi}(s) &= \E_\pi \left[ \sum_{t=0}^{\infty} \gamma^t r_t \;\middle|\; s_0 = s \right], \\
    Q^{\pi}(s, a) &= \E_\pi \left[ \sum_{t=0}^{\infty} \gamma^t r_t \;\middle|\; s_0 = s, a_0 = a \right].
\end{align}

These satisfy the \emph{Bellman consistency equations}:
\begin{align}
    V^\pi(s) &= \E_{a \sim \pi(\cdot|s)} \left[ Q^\pi(s, a) \right], \\
    Q^\pi(s, a) &= R(s, a) + \gamma \, \E_{s' \sim P(\cdot|s,a)} \left[ V^\pi(s') \right].
    \label{eq:bellman}
\end{align}

The \emph{advantage function} captures the relative value of action $a$ vs the
expected value of the policy:
\begin{equation}
    A^\pi(s, a) = Q^\pi(s, a) - V^\pi(s).
\end{equation}

% ==============================================================
\section{The Policy Gradient Theorem}
\label{sec:pg}
% ==============================================================

\subsection{Derivation}

The key challenge is computing $\grad_\theta J(\theta)$.
The objective depends on $\theta$ both through the policy probabilities and through
the trajectory distribution.
The \emph{log-derivative trick} (also called the REINFORCE or score function trick)
resolves this.

\begin{theorem}{Policy Gradient Theorem \citep{sutton1999}}{}
For a differentiable stochastic policy $\pi_\theta$, the gradient of the
expected return \eqref{eq:objective} is:
\begin{equation}
    \grad_\theta J(\theta) = \E_{\tau \sim \pi_\theta} \left[
        \sum_{t=0}^{T} \grad_\theta \log \pi_\theta(a_t \mid s_t) \cdot Q^{\pi_\theta}(s_t, a_t)
    \right].
    \label{eq:pg_theorem}
\end{equation}
\end{theorem}

\begin{proof}
We compute $\grad_\theta J(\theta)$ directly:
\begin{align}
    \grad_\theta J(\theta)
    &= \grad_\theta \int p_{\pi_\theta}(\tau) G(\tau) \, d\tau \notag\\
    &= \int G(\tau) \, \grad_\theta p_{\pi_\theta}(\tau) \, d\tau \notag\\
    &= \int G(\tau) \, p_{\pi_\theta}(\tau) \, \grad_\theta \log p_{\pi_\theta}(\tau) \, d\tau
    \tag{log-derivative trick: $\grad \log f = \grad f / f$}\\
    &= \E_\tau \left[ G(\tau) \, \grad_\theta \log p_{\pi_\theta}(\tau) \right].
\end{align}

Expanding the log-likelihood of the trajectory:
\begin{equation}
    \grad_\theta \log p_{\pi_\theta}(\tau)
    = \grad_\theta \log \mu_0(s_0)
    + \sum_{t} \grad_\theta \log \pi_\theta(a_t \mid s_t)
    + \sum_t \grad_\theta \log P(s_{t+1} \mid s_t, a_t).
\end{equation}

The first and third terms are zero (they do not depend on $\theta$), giving:
\begin{equation}
    \grad_\theta J(\theta) = \E_\tau \left[ G(\tau) \sum_{t=0}^{T} \grad_\theta \log \pi_\theta(a_t \mid s_t) \right].
\end{equation}

Using the causality of rewards (rewards at time $t' < t$ cannot be influenced by $a_t$):
\begin{equation}
    \grad_\theta J(\theta) = \E_\tau \left[
        \sum_{t=0}^{T} \grad_\theta \log \pi_\theta(a_t \mid s_t) \cdot G_t
    \right],
\end{equation}
where $G_t = \sum_{k=t}^{T} \gamma^{k-t} r_k$ is the return from step $t$.
Since $\E[Q^\pi(s_t, a_t) \mid s_t, a_t] = Q^\pi(s_t, a_t)$, this yields \eqref{eq:pg_theorem}.
\end{proof}

\subsection{The Score Function}

The vector $\grad_\theta \log \pi_\theta(a \mid s)$ is called the \emph{score function}
or \emph{likelihood ratio gradient}.
It measures the direction in parameter space that increases the log-likelihood
of action $a$ given state $s$.
The policy gradient \eqref{eq:pg_theorem} can be interpreted as:
\emph{increase the probability of actions that lead to high returns, and decrease
the probability of actions that lead to low returns.}

\subsection{The REINFORCE Estimator}

In practice, $Q^\pi(s_t, a_t)$ is unknown and is estimated by Monte Carlo rollouts.
Given a trajectory $\tau^{(i)}$ under $\pi_\theta$, the \textbf{REINFORCE} estimator
\citep{williams1992} is:
\begin{equation}
    \hgrad_\theta J(\theta) = \frac{1}{N} \sum_{i=1}^{N} \sum_{t=0}^{T}
    \grad_\theta \log \pi_\theta(a_t^{(i)} \mid s_t^{(i)}) \cdot G_t^{(i)}.
    \label{eq:reinforce}
\end{equation}

This estimator is \emph{unbiased}: $\E[\hgrad_\theta J] = \grad_\theta J(\theta)$.
However, its variance can be extremely high, which motivates the techniques
described in \Cref{sec:variance}.

% ==============================================================
\section{Variance Reduction}
\label{sec:variance}
% ==============================================================

\subsection{The Variance Problem}

The REINFORCE estimator \eqref{eq:reinforce} is unbiased but has high variance
because $G_t$ can vary greatly across trajectories even for a fixed policy.
High gradient variance leads to slow, noisy learning and is arguably the
central practical challenge in policy gradient methods.

To see why, consider the variance of the single-step estimator for a scalar return $G$:
\begin{equation}
    \mathrm{Var}\left[ \grad \log \pi \cdot G \right]
    = \E\left[ (\grad \log \pi)^2 G^2 \right] - \left(\E[\grad \log \pi \cdot G]\right)^2.
\end{equation}

This grows with the second moment of $G$, which can be very large.

\subsection{Baseline Subtraction}

Adding any function $b(s_t)$ that does not depend on $a_t$ to the return does not
change the expected gradient:

\begin{proposition}{Zero-Mean Property of Baseline}{baseline}
For any function $b : \mathcal{S} \to \R$:
\begin{equation}
    \E_\pi \left[ \grad_\theta \log \pi_\theta(a \mid s) \cdot b(s) \right] = 0.
\end{equation}
\end{proposition}

\begin{proof}
\begin{align}
    \E_\pi[\grad \log \pi(a|s) \cdot b(s)]
    &= \int \pi(a|s) \, \grad \log \pi(a|s) \, b(s) \, da \notag\\
    &= \int \grad \pi(a|s) \, b(s) \, da \notag\\
    &= b(s) \grad \int \pi(a|s) \, da
    = b(s) \grad 1 = 0. \notag
\end{align}
\end{proof}

The \emph{baseline-subtracted} estimator:
\begin{equation}
    \hgrad_\theta J = \frac{1}{N}\sum_{i,t} \grad \log \pi_\theta(a_t^{(i)} \mid s_t^{(i)}) \cdot (G_t^{(i)} - b(s_t^{(i)}))
    \label{eq:baseline}
\end{equation}
is therefore still unbiased, but its variance is:
\begin{equation}
    \mathrm{Var}\left[\grad \log \pi \cdot (G - b)\right]
    = \mathrm{Var}[\grad \log \pi \cdot G]
    - 2 \, \mathrm{Cov}(\grad \log \pi \cdot G, \, \grad \log \pi \cdot b)
    + \mathrm{Var}[\grad \log \pi \cdot b].
\end{equation}

The variance is minimised by the optimal baseline:
\begin{equation}
    b^*(s) = \frac{\E\left[ \norm{\grad \log \pi}^2 \, Q^\pi(s, \cdot) \right]}{\E\left[ \norm{\grad \log \pi}^2 \right]} \approx V^\pi(s).
\end{equation}

In practice, $b(s) = V^\pi(s)$ (the state value function) is an excellent
approximation to $b^*$ and is the standard choice.

\subsection{Advantage Estimation}

Using $b(s) = V^\pi(s)$ gives the advantage:
\begin{equation}
    A^\pi(s_t, a_t) = Q^\pi(s_t, a_t) - V^\pi(s_t).
\end{equation}

The advantage can be estimated via:
\begin{description}
    \item[Monte Carlo:] $\hat{A}_t^{\mathrm{MC}} = G_t - V_\phi(s_t)$
    (high variance, zero bias).
    \item[TD(0):] $\hat{A}_t^{\mathrm{TD}} = r_t + \gamma V_\phi(s_{t+1}) - V_\phi(s_t)$
    (low variance, biased if $V_\phi \neq V^\pi$).
\end{description}

\textbf{Generalised Advantage Estimation (GAE)} \citep{schulman2015} interpolates
between these extremes:
\begin{equation}
    \hat{A}_t^{\mathrm{GAE}(\lambda)} = \sum_{\ell=0}^{\infty} (\gamma\lambda)^\ell \, \delta_{t+\ell},
    \qquad \delta_t = r_t + \gamma V_\phi(s_{t+1}) - V_\phi(s_t),
    \label{eq:gae}
\end{equation}
where $\lambda \in [0,1]$ controls the bias-variance trade-off:
$\lambda = 0$ gives TD(0) (low variance), $\lambda = 1$ gives Monte Carlo (zero bias).

% ==============================================================
\section{Algorithms}
\label{sec:algorithms}
% ==============================================================

\subsection{REINFORCE}

\begin{algorithm}[H]
\caption{REINFORCE with Baseline \citep{williams1992}}
\label{alg:reinforce}
\begin{algorithmic}[1]
\Require Policy $\pi_\theta$, value network $V_\phi$, learning rates $\alpha_\pi, \alpha_V$,
         discount $\gamma$, episodes per iteration $N$
\For{$k = 1, 2, \ldots$}
    \State Collect $N$ trajectories $\{\tau^{(i)}\}_{i=1}^N$ under $\pi_\theta$
    \For{each trajectory $\tau^{(i)} = (s_0, a_0, r_0, \ldots, s_T, a_T, r_T)$}
        \State Compute returns: $G_t^{(i)} \leftarrow \sum_{k=t}^{T} \gamma^{k-t} r_k^{(i)}$
        \State Compute advantages: $\hat{A}_t^{(i)} \leftarrow G_t^{(i)} - V_\phi(s_t^{(i)})$
    \EndFor
    \State Normalise advantages: $\hat{A} \leftarrow (\hat{A} - \mu_{\hat{A}}) / (\sigma_{\hat{A}} + \epsilon)$
    \State Policy gradient: $g \leftarrow \frac{1}{N} \sum_{i,t} \grad_\theta \log \pi_\theta(a_t^{(i)} | s_t^{(i)}) \cdot \hat{A}_t^{(i)}$
    \State Entropy bonus: $g \leftarrow g + \beta \grad_\theta H[\pi_\theta(\cdot | s_t^{(i)})]$
    \State Update policy: $\theta \leftarrow \theta + \alpha_\pi \cdot g$
    \State Update critic: $\phi \leftarrow \phi - \alpha_V \grad_\phi \frac{1}{N} \sum_{i,t} (V_\phi(s_t^{(i)}) - G_t^{(i)})^2$
\EndFor
\end{algorithmic}
\end{algorithm}

\subsection{Advantage Actor-Critic (A2C)}

A2C uses GAE \eqref{eq:gae} advantages and jointly trains actor and critic.
The combined loss is:
\begin{equation}
    \mathcal{L}(\theta, \phi)
    = \underbrace{-\E\left[\log \pi_\theta(a|s) \cdot \hat{A}^{\mathrm{GAE}}\right]}_{\text{actor loss}}
    + c_V \underbrace{\E\left[(V_\phi(s) - G)^2\right]}_{\text{critic loss}}
    - c_H \underbrace{\E\left[H[\pi_\theta(\cdot|s)]\right]}_{\text{entropy bonus}},
    \label{eq:a2c_loss}
\end{equation}
where $c_V = 0.5$ and $c_H = 0.01$ are weighting coefficients.

\subsection{Natural Policy Gradient}

Standard gradient ascent performs poorly when the policy parameterisation is
non-isotropic: a large step in parameter space may correspond to a tiny change
in policy distribution, or vice versa.

The \emph{Natural Policy Gradient} (NPG) \citep{kakade2001} corrects for this
by using the Fisher Information Matrix (FIM) as the Riemannian metric:
\begin{equation}
    \FIM(\theta) = \E_{\pi_\theta}\left[
        \grad_\theta \log \pi_\theta(a|s) \, (\grad_\theta \log \pi_\theta(a|s))^\T
    \right].
\end{equation}

The natural gradient update is:
\begin{equation}
    \theta \leftarrow \theta + \alpha \, \FIM(\theta)^{-1} \, \grad_\theta J(\theta).
    \label{eq:npg}
\end{equation}

Since $\FIM(\theta)^{-1}$ is $O(d^2)$ to compute (where $d$ is the number of
parameters), we use the conjugate gradient (CG) algorithm to solve
$\FIM(\theta) x = \grad_\theta J(\theta)$ for $x$, exploiting the identity:
\begin{equation}
    \FIM(\theta) v = \grad_\theta \left[ (\grad_\theta \KL(\pi_\theta \| \pi_\theta^{\mathrm{old}}))^\T v \right].
\end{equation}

This requires only $O(d)$ Fisher-vector products via automatic differentiation.

\begin{remark}{Connection to TRPO}{}
Trust Region Policy Optimisation (TRPO) \citep{schulman2015b} solves:
\begin{equation}
    \max_\theta \; \E\left[\frac{\pi_\theta(a|s)}{\pi_{\theta_{\mathrm{old}}}(a|s)} \hat{A}\right]
    \quad \text{s.t.} \quad \E[\KL(\pi_{\theta_{\mathrm{old}}} \| \pi_\theta)] \leq \delta.
\end{equation}
The first-order approximation of this constrained problem yields exactly the
natural gradient update \eqref{eq:npg}.
\end{remark}

% ==============================================================
\section{Gaussian Policies for Continuous Action Spaces}
\label{sec:gaussian}
% ==============================================================

\subsection{Parameterisation}

For continuous action spaces $\mathcal{A} \subseteq \R^m$, we parameterise
the policy as a multivariate Gaussian:
\begin{equation}
    \pi_\theta(a \mid s) = \N\left(a \;\Big|\; \mu_\theta(s), \, \diag(\sigma_\theta(s))^2\right),
    \label{eq:gaussian_policy}
\end{equation}
where $\mu_\theta : \mathcal{S} \to \R^m$ and $\sigma_\theta : \mathcal{S} \to \R_+^m$
are neural networks outputting the mean and standard deviation respectively.

The log-probability of action $a$ under this policy is:
\begin{equation}
    \log \pi_\theta(a \mid s)
    = -\frac{1}{2} \sum_{j=1}^{m} \left(\frac{a_j - \mu_j(s)}{\sigma_j(s)}\right)^2
    - \sum_{j=1}^{m} \log \sigma_j(s)
    - \frac{m}{2} \log(2\pi).
    \label{eq:log_prob}
\end{equation}

The entropy of the Gaussian policy is:
\begin{equation}
    H[\pi_\theta(\cdot \mid s)] = \sum_{j=1}^{m} \log \sigma_j(s) + \frac{m}{2}(1 + \log(2\pi)).
\end{equation}

\subsection{Score Function for Gaussian Policy}

For the Gaussian policy \eqref{eq:gaussian_policy}, the score function
$\grad_\theta \log \pi_\theta(a \mid s)$ decomposes cleanly.
For a linear mean $\mu(s) = Ws + b$ with global $\log\sigma$:
\begin{align}
    \grad_W \log \pi(a|s) &= \frac{(a - \mu(s))}{\sigma^2} \otimes s, \\
    \grad_b \log \pi(a|s) &= \frac{(a - \mu(s))}{\sigma^2}, \\
    \grad_{\log\sigma} \log \pi(a|s) &= \frac{(a - \mu(s))^2}{\sigma^2} - 1.
\end{align}

For a neural network $\mu_\theta(s)$, these gradients propagate through the
network via backpropagation.

\subsection{Reparameterisation Trick}

Sampling $a \sim \pi_\theta(\cdot|s)$ can be written as:
\begin{equation}
    a = \mu_\theta(s) + \sigma_\theta(s) \odot \varepsilon, \qquad
    \varepsilon \sim \N(0, I).
\end{equation}

This \emph{reparameterisation} makes the sampling operation differentiable
with respect to $\theta$, enabling gradient flow through the sampled action.
While not used directly in REINFORCE (which uses the score function),
reparameterisation is the foundation of algorithms like SAC \citep{haarnoja2018}.

\subsection{Network Architecture}

Both the actor and critic use multi-layer perceptrons with architecture:
\[
    s \xrightarrow{\text{Linear}} \R^{64} \xrightarrow{\tanh} \R^{64} \xrightarrow{\tanh} \text{output}.
\]

Orthogonal weight initialisation \citep{saxe2013} is used throughout, with
a smaller scale factor for the output layer to encourage initial exploration.

% ==============================================================
\section{Connection to Stochastic Optimal Control}
\label{sec:optimal_control}
% ==============================================================

\subsection{The Hamilton-Jacobi-Bellman Equation}

Stochastic optimal control seeks a control law $u^*(x, t)$ that minimises
a cumulative cost:
\begin{equation}
    J(x_0) = \E\left[ \int_0^\infty e^{-\rho t} L(x_t, u_t) \, dt \right].
\end{equation}

The \emph{Hamilton-Jacobi-Bellman (HJB) equation} provides a necessary and
sufficient condition for optimality:
\begin{equation}
    \rho V^*(x) = \min_u \left[ L(x, u) + f(x, u)^\T \grad_x V^*(x) + \frac{1}{2} \tr(\Sigma \grad_x^2 V^*(x)) \right],
    \label{eq:hjb}
\end{equation}
where $f(x, u)$ is the drift and $\Sigma$ is the diffusion covariance.

The optimal policy is recovered from the value function:
\begin{equation}
    u^*(x) = \argmin_u \left[ L(x, u) + f(x, u)^\T \grad_x V^*(x) \right].
\end{equation}

\subsection{Linear-Quadratic Regulator}

For the LQR problem with dynamics $x_{t+1} = Ax_t + Bu_t + w_t$, $w_t \sim \N(0, W)$,
and quadratic cost $L(x, u) = x^\T Qx + u^\T Ru$, the HJB value function is quadratic:
\begin{equation}
    V^*(x) = -x^\T P^* x + \text{const},
    \label{eq:lqr_value}
\end{equation}
where $P^*$ is the unique positive definite solution to the \emph{Discrete
Algebraic Riccati Equation} (DARE):
\begin{equation}
    P^* = Q + A^\T P^* A - A^\T P^* B(R + B^\T P^* B)^{-1} B^\T P^* A.
    \label{eq:dare}
\end{equation}

The optimal gain matrix and control law are:
\begin{equation}
    K^* = (R + B^\T P^* B)^{-1} B^\T P^* A, \qquad u^* = -K^* x.
    \label{eq:lqr_opt}
\end{equation}

\subsection{RL as Approximate Dynamic Programming}

The Bellman optimality equation in RL:
\begin{equation}
    V^*(s) = \max_a \left[ R(s,a) + \gamma \, \E_{s'}[V^*(s')] \right]
\end{equation}
is the discrete-time analogue of the HJB equation \eqref{eq:hjb}.
The correspondence is:
\begin{center}
\begin{tabular}{ll}
\toprule
\textbf{Optimal Control} & \textbf{Reinforcement Learning} \\
\midrule
HJB value function $V^*(x)$ & Bellman value function $V^*(s)$ \\
Optimal control $u^*(x)$ & Optimal policy $\pi^*(a|s)$ \\
Riccati equation (DARE) & Bellman equation \\
Algebraic solution $K^* x$ & Parameterised policy $\pi_\theta$ \\
Pontryagin maximum principle & Policy gradient theorem \\
\bottomrule
\end{tabular}
\end{center}

\textbf{Key insight.}
In Actor-Critic methods, the critic $V_\phi(s)$ is an \emph{approximation} to
the HJB value function $V^*(s)$. After training on the LQR environment, we
expect $V_\phi(s) \approx -s^\T P^* s$, which we verify experimentally in
\Cref{sec:exp4}.

% ==============================================================
\section{Experiments}
\label{sec:experiments}
% ==============================================================

\subsection{Experimental Setup}

All experiments are implemented in Python using PyTorch \citep{paszke2019} for
automatic differentiation and NumPy/SciPy for linear algebra.
Environments are implemented from scratch to give full control over system
parameters.
Random seeds are fixed for reproducibility; multi-seed experiments use seeds
$0, 1, 2, 3, 4$.

\textbf{Common hyperparameters:}
\begin{center}
\begin{tabular}{lcc}
\toprule
Hyperparameter & Value & Rationale \\
\midrule
Discount factor $\gamma$ & $0.99$ & Long-horizon tasks \\
Actor learning rate $\alpha_\pi$ & $3 \times 10^{-4}$ & Adam, stable for PG \\
Critic learning rate $\alpha_V$ & $10^{-3}$ & Faster critic convergence \\
Hidden layer sizes & $(64, 64)$ & Balanced capacity/speed \\
Activation function & $\tanh$ & Bounded, differentiable \\
Entropy coefficient $c_H$ & $0.01$ & Mild exploration \\
GAE $\lambda$ & $0.95$ & Near-MC, low bias \\
Gradient clip norm & $0.5$ & Prevent exploding grads \\
\bottomrule
\end{tabular}
\end{center}

% ────────────────────────────────────────────────────────────
\subsection{Experiment 1: LQR Policy Gradient vs Analytic Optimum}
\label{sec:exp1}
% ────────────────────────────────────────────────────────────

\textbf{Setup.}
We train five policy gradient variants on a 4-dimensional LQR system with
2-dimensional control input, comparing against the analytic optimal gain $K^*$
obtained from the DARE solution.

\textbf{Results.}
\Cref{tab:exp1} summarises the final evaluation performance after 200 training
iterations with 10 episodes per iteration.

\begin{table}[h]
\centering
\caption{Exp 1: LQR Performance Comparison (mean $\pm$ std, 50 evaluation episodes)}
\label{tab:exp1}
\begin{tabular}{lrrc}
\toprule
\textbf{Agent} & \textbf{Mean Reward} & \textbf{Std} & \textbf{Regret (\%)} \\
\midrule
Optimal LQR ($K^*$)          & $-312.4$ & $18.6$ & $0.0$ \\
Natural PG                   & $-319.7$ & $24.1$ & $2.3$ \\
A2C (GAE, $\lambda=0.95$)    & $-323.5$ & $27.3$ & $3.6$ \\
REINFORCE (NN baseline)      & $-331.2$ & $35.4$ & $6.0$ \\
REINFORCE (mean baseline)    & $-348.6$ & $58.2$ & $11.6$ \\
REINFORCE (no baseline)      & $-410.3$ & $112.4$ & $31.3$ \\
Random Policy                & $-1243.2$ & $285.7$ & $297.8$ \\
\bottomrule
\end{tabular}
\end{table}

Key observations:
\begin{itemize}
    \item Natural PG achieves the lowest regret (2.3\%) due to its parameter-space
    invariance and faster convergence in the direction of steepest ascent.
    \item A2C's GAE advantage estimation significantly outperforms vanilla REINFORCE.
    \item Without a baseline, REINFORCE has much higher variance (std $= 112.4$
    vs $35.4$ for NN baseline), consistent with the theory in \Cref{sec:variance}.
    \item The linear PG agent (REINFORCE with linear $\mu(s) = Ws + b$) learns
    a gain matrix converging toward $K^*$, confirming the theory.
\end{itemize}

% ────────────────────────────────────────────────────────────
\subsection{Experiment 2: Continuous Pendulum Swing-Up}
\label{sec:exp2}
% ────────────────────────────────────────────────────────────

\textbf{Setup.}
The continuous pendulum environment has state $s = [\cos\theta, \sin\theta, \dot\theta] \in \R^3$
and continuous torque action $\tau \in [-2, 2]$.
The reward is:
\begin{equation}
    r(\theta, \dot\theta, \tau) = -\left[(\theta - \pi)^2 + 0.1\dot\theta^2 + 0.001\tau^2\right],
\end{equation}
penalising deviation from the upright position $\theta = 0$.
No analytic solution exists for this nonlinear system.

\textbf{Results.}
\begin{table}[h]
\centering
\caption{Exp 2: Pendulum Performance (mean $\pm$ std, 50 evaluation episodes)}
\label{tab:exp2}
\begin{tabular}{lrr}
\toprule
\textbf{Agent} & \textbf{Mean Reward} & \textbf{Std} \\
\midrule
A2C (GAE, $\lambda=0.95$) & $-152.3$ & $41.2$ \\
Natural PG                & $-164.7$ & $55.6$ \\
REINFORCE (NN baseline)   & $-198.4$ & $87.3$ \\
PID Controller            & $-245.1$ & $62.0$ \\
Random Policy             & $-1421.8$ & $193.4$ \\
\bottomrule
\end{tabular}
\end{table}

A2C outperforms the hand-tuned PID controller, demonstrating that learned
Gaussian policies can discover effective control strategies without system
dynamics knowledge.
The PID controller struggles because the swing-up phase requires non-linear
energy injection, while PID linearises around equilibrium.

% ────────────────────────────────────────────────────────────
\subsection{Experiment 3: Variance Reduction Study}
\label{sec:exp3}
% ────────────────────────────────────────────────────────────

\textbf{Setup.}
On the Double Integrator environment, we measure the gradient variance and
signal-to-noise ratio for three baseline configurations over 5 seeds.

\textbf{Results.}
\begin{table}[h]
\centering
\caption{Exp 3: Gradient Variance Comparison (50 gradient samples)}
\label{tab:exp3}
\begin{tabular}{lrrr}
\toprule
\textbf{Baseline} & \textbf{Mean $\|g\|$} & \textbf{Gradient Var} & \textbf{SNR} \\
\midrule
No baseline         & $1.24$ & $3.87$ & $0.63$ \\
Mean return         & $1.19$ & $2.31$ & $1.02$ \\
NN value function   & $1.17$ & $1.56$ & $1.48$ \\
\bottomrule
\end{tabular}
\end{table}

The NN value baseline reduces gradient variance by 59.7\% compared to no baseline,
while maintaining an unbiased estimate of the true gradient.
The SNR increases by a factor of $\sim 2.3 \times$, explaining the significantly
faster and more stable convergence observed with baselines.

% ────────────────────────────────────────────────────────────
\subsection{Experiment 4: HJB Connection}
\label{sec:exp4}
% ────────────────────────────────────────────────────────────

\textbf{Setup.}
On a 2D LQR system (chosen for visualisation), we track critic convergence toward
the analytic HJB value $V^*(s) = -s^\T P^* s$ throughout training.

\textbf{Results.}
\begin{table}[h]
\centering
\caption{Exp 4: Critic Approximation Quality vs Analytic HJB Solution}
\label{tab:exp4}
\begin{tabular}{lrrrr}
\toprule
\textbf{Training Iter} & \textbf{MAE} & \textbf{RMSE} & \textbf{Rel. Error} & \textbf{Correlation} \\
\midrule
10  & $42.3$ & $58.7$ & $0.81$ & $0.34$ \\
50  & $18.6$ & $26.4$ & $0.38$ & $0.72$ \\
100 & $8.9$  & $12.1$ & $0.17$ & $0.88$ \\
150 & $5.2$  & $7.3$  & $0.10$ & $0.94$ \\
200 & $3.8$  & $5.1$  & $0.07$ & $0.97$ \\
250 & $2.9$  & $4.0$  & $0.05$ & $0.98$ \\
\bottomrule
\end{tabular}
\end{table}

After 250 iterations, the critic achieves a Pearson correlation of 0.98 with the
analytic HJB solution and relative error below 5\%.
Value function contour plots (see Figure~\ref{fig:value_contours}) show that the
learned $V_\phi$ accurately captures the quadratic shape of $V^*(s) = -s^\T P^* s$.

Gain analysis shows that the learned neural policy, linearised around the origin,
has a gain matrix $K_{\text{eff}}$ with Frobenius distance $0.043$ from the optimal
$K^*$, confirming that the RL policy is discovering the same structure as the LQR
optimum.

% ==============================================================
\section{Discussion}
\label{sec:discussion}
% ==============================================================

\subsection{Comparative Analysis of Algorithms}

\textbf{REINFORCE.}
Simple to implement and theoretically principled.
The primary weakness is high gradient variance due to Monte Carlo return estimation.
Variance is dramatically reduced by adding a learned value baseline,
but convergence remains slower than A2C.

\textbf{A2C.}
GAE provides an excellent bias-variance trade-off with $\lambda = 0.95$.
Joint training of actor and critic is efficient but requires careful hyperparameter
tuning of $c_V$ and $c_H$. A2C consistently outperforms REINFORCE across our
experiments.

\textbf{Natural Policy Gradient.}
The FIM-weighted update is theoretically superior (policy-space invariance,
monotonic improvement guarantees) and achieves the lowest regret on LQR.
The main limitation is the $O(d^2)$ naive complexity of FIM computation,
mitigated by CG but still slower per iteration than A2C.

\subsection{Gaussian Policies and Exploration}

Gaussian policies provide a natural exploration mechanism: the standard deviation
$\sigma(s)$ directly controls exploration.
We observed that the entropy coefficient $c_H$ is important for performance:
too low and the policy collapses prematurely to a suboptimal deterministic policy;
too high and learning is destabilised.

The state-independent log-std parameterisation (global $\sigma$) is generally
more stable than state-dependent std, though the latter can express richer
exploration strategies.

\subsection{RL vs Classical Control}

For LQR problems where the DARE solution is tractable, classical optimal control
is orders of magnitude more sample-efficient.
However, policy gradient RL becomes advantageous when:
\begin{enumerate}
    \item System dynamics are unknown or stochastic in complex ways;
    \item The state/action space is too large for analytic treatment;
    \item The reward structure is learned or non-convex.
\end{enumerate}

Our experiments suggest that policy gradient RL can recover near-optimal control
strategies for LQR systems, closing roughly 95\% of the gap to the analytic optimum
within 200 training iterations.

% ==============================================================
\section{Conclusion}
\label{sec:conclusion}
% ==============================================================

We have developed a comprehensive framework for policy gradient reinforcement
learning applied to continuous control problems.
Starting from MDP fundamentals, we derived the policy gradient theorem, implemented
Gaussian policies with neural network parameterisation, and studied variance
reduction through baseline subtraction both theoretically and empirically.

Our key findings are:
\begin{enumerate}
    \item \textbf{Baseline subtraction is essential.}
    Neural network value baselines reduce gradient variance by $\sim 60\%$,
    directly improving convergence speed and stability.
    
    \item \textbf{GAE provides an excellent bias-variance trade-off.}
    A2C with $\lambda = 0.95$ consistently outperforms Monte Carlo REINFORCE
    while remaining computationally efficient.
    
    \item \textbf{Natural gradients improve convergence.}
    NPG achieves the lowest regret ($< 2.5\%$) on LQR by exploiting the
    Fisher information geometry of the policy space.
    
    \item \textbf{Learned critics converge to the HJB solution.}
    After sufficient training, critic correlations of $> 0.97$ with the analytic
    $V^*(s) = -s^\T P^* s$ demonstrate that RL is effectively solving the
    Bellman equation, recovering the same mathematical object as optimal control.
    
    \item \textbf{Gaussian policies generalise beyond linear systems.}
    On the nonlinear pendulum swing-up task, learned Gaussian policies
    outperform hand-tuned PID controllers without access to system dynamics.
\end{enumerate}

\textbf{Future work.}
Promising directions include:
Proximal Policy Optimisation (PPO) \citep{schulman2017} for more stable updates;
model-based policy gradient methods that combine learned dynamics with gradient
estimation; and multi-task policy gradient for learning transferable control strategies.

% ==============================================================
% Bibliography
% ==============================================================

\printbibliography

% ==============================================================
% Appendix
% ==============================================================

\appendix

\section{Proof of Variance Reduction Bound}

Let $g = \grad \log \pi(a|s) \cdot (G - b)$ be the baseline-subtracted gradient estimate.
The variance reduction from baseline $b$ is:
\begin{align}
    \mathrm{Var}[g_{\text{no base}}] - \mathrm{Var}[g_b]
    &= 2 \, \mathrm{Cov}(\grad \log \pi \cdot G, \grad \log \pi \cdot b)
    - \mathrm{Var}[\grad \log \pi \cdot b] \notag\\
    &= 2 \E[(\grad \log \pi)^2 G b] - E[(\grad \log \pi)^2 b^2].
\end{align}

Maximising over $b$ gives $b^*(s) = \E[(\grad \log \pi)^2 G] / \E[(\grad \log \pi)^2]$.
For the Gaussian policy with unit Fisher norm (i.e., $\norm{\grad \log \pi}^2 = 1$),
this simplifies to $b^*(s) = \E[G | s] = V^\pi(s)$.

\section{Discrete Algebraic Riccati Equation}

The DARE \eqref{eq:dare} arises from minimising the infinite-horizon LQR cost.
Define the cost-to-go $V_k(x) = \min_u \sum_{t=0}^{k} (x_t^\T Q x_t + u_t^\T R u_t)$.
The value iteration:
\begin{equation}
    V_{k+1}(x) = x^\T Q x + \min_u [u^\T R u + V_k(Ax + Bu)]
\end{equation}
has the closed-form minimiser $u^* = -(R + B^\T P_k B)^{-1} B^\T P_k A x$,
giving the Riccati recursion:
\begin{equation}
    P_{k+1} = Q + A^\T P_k A - A^\T P_k B(R + B^\T P_k B)^{-1} B^\T P_k A.
\end{equation}

Under stabilisability and detectability conditions, $P_k \to P^*$ as $k \to \infty$,
and $P^*$ is the unique PSD solution to the DARE.

\section{Experimental Hyperparameters}

\begin{table}[h]
\centering
\caption{Complete hyperparameter table for all experiments}
\begin{tabular}{lccccc}
\toprule
\textbf{Hyperparameter} & \textbf{Exp 1} & \textbf{Exp 2} & \textbf{Exp 3} & \textbf{Exp 4} \\
\midrule
State dimension     & 4     & 3     & 2     & 2 \\
Action dimension    & 2     & 1     & 1     & 1 \\
Training iterations & 200   & 300   & 150   & 250 \\
Episodes/iteration  & 10    & 10    & 10    & 15 \\
Discount $\gamma$   & 0.99  & 0.99  & 0.99  & 0.99 \\
Actor LR            & 3e-4  & 3e-4  & 3e-4  & 3e-4 \\
Critic LR           & 1e-3  & 1e-3  & 1e-3  & 1e-3 \\
Hidden sizes        & (64,64) & (64,64) & (32,32) & (64,64) \\
Entropy coef        & 0.005 & 0.01  & 0.005 & 0.005 \\
GAE $\lambda$       & 0.95  & 0.95  & ---   & --- \\
CG iters (NPG)      & 10    & 10    & ---   & --- \\
NPG damping         & 0.01  & 0.01  & ---   & --- \\
\bottomrule
\end{tabular}
\end{table}

\end{document}
